\documentclass[UTF8,fontset=none,11pt,a4paper]{ctexart}
\usepackage[left=21mm,right=21mm,top=27mm,bottom=22mm,headheight=26pt,headsep=6mm]{geometry}
\usepackage{fontspec}
\setmainfont{texgyretermes-regular.otf}[BoldFont=texgyretermes-bold.otf,ItalicFont=texgyretermes-italic.otf]
\setCJKmainfont{Songti.ttc}[Path=/System/Library/Fonts/Supplemental/,FontIndex=6,AutoFakeBold=2]
\setCJKsansfont{Songti.ttc}[Path=/System/Library/Fonts/Supplemental/,FontIndex=6,AutoFakeBold=2]
\usepackage[table]{xcolor}
\definecolor{TableHead}{HTML}{EEF0ED}
\definecolor{Muted}{HTML}{666C68}
\usepackage{booktabs,longtable,array,fancyhdr,needspace,xurl}
\usepackage[unicode,hidelinks]{hyperref}
\urlstyle{same}
\linespread{1.32}
\setlength{\parindent}{2em}
\setlength{\parskip}{3pt}
\setlength{\LTpre}{8pt}\setlength{\LTpost}{10pt}
\ctexset{section={format=\fontsize{15}{20}\selectfont\bfseries,beforeskip=18pt,afterskip=8pt},subsection={format=\fontsize{12}{16}\selectfont\bfseries,beforeskip=12pt,afterskip=6pt},subsubsection={format=\normalsize\bfseries,beforeskip=9pt,afterskip=4pt}}
\pagestyle{fancy}\fancyhf{}
\renewcommand{\headrulewidth}{.3pt}
\renewcommand{\footrulewidth}{0pt}
\fancyhead[L]{\parbox[b]{\linewidth}{\fontsize{8}{10}\selectfont\color{Muted}教育贷款暂停还款：资格与利息安排·本次账户资料及办理记录}}
\fancyfoot[C]{\small\color{Muted}\thepage}
\hypersetup{pdftitle={本次账户资料及办理记录},pdfsubject={仅供模拟使用·教育贷款暂停还款：资格与利息安排·14-03},pdfauthor={金融决策模拟材料}}
\begin{document}
{\noindent\fontsize{20}{25}\selectfont\bfseries 本次账户资料及办理记录\par}
\vspace{7pt}{\noindent\small\color{Muted}仅供模拟使用·教育贷款暂停还款：资格与利息安排·14-03\par}\vspace{10pt}
\subsection*{A：一般暂缓还款}
\textbf{所用机制：} General Forbearance


\textbf{申请依据：} 当前暂时性经济困难，服务机构审核批准


\textbf{本次暂停覆盖：} 两笔贷款，同一个90天期间


\textbf{期间正常月供：} 暂停要求支付


\textbf{补贴贷款本金：} 4,500美元


\textbf{补贴贷款期间利息：} 按4.53\%计入借款人应付利息


\textbf{非补贴贷款本金：} 5,500美元


\textbf{非补贴贷款期间利息：} 按4.53\%计入借款人应付利息


\textbf{申请资料：} 暂缓申请表，说明暂时性经济困难及所需期间


\textbf{本次资料准备状态：} 可以完整提交


\textbf{原贷款合同利率：} 保持不变


\textbf{提交及批准费用：} 无本次服务费


\textbf{结束处理：} 恢复付款；本次安排支付期间新增利息


\subsection*{B：失业延期还款}
\textbf{所用机制：} Unemployment Deferment


\textbf{申请依据：} 当前失业及相应补助证明，核对延期资格


\textbf{本次暂停覆盖：} 两笔贷款，同一个90天期间


\textbf{期间正常月供：} 暂停要求支付


\textbf{补贴贷款本金：} 4,500美元


\textbf{补贴贷款期间利息：} 该延期期间不向借款人收取利息


\textbf{非补贴贷款本金：} 5,500美元


\textbf{非补贴贷款期间利息：} 同样按4.53\%计入借款人应付利息


\textbf{申请资料：} 延期申请表，并附失业补助证明


\textbf{本次资料准备状态：} 正式补助证明已取得，可以完整提交


\textbf{原贷款合同利率：} 保持不变


\textbf{提交及批准费用：} 无本次服务费


\textbf{结束处理：} 恢复付款；本次安排支付期间新增利息


\subsection*{办理确认事项}
\textbf{资格与期限。} B使用已取得的失业补助证明申请，不需要在本次任务中另找资格证据。两种机制均有可用期限，且服务机构已确认90天能够覆盖；不会因一边申请不获批或批准较晚而缩短暂停期。


\textbf{利息计算对象。} 两笔贷款保持各自的本金和贷款类型，未合并或转为私人贷款。期间没有本金还款或新的借款，利息按照各笔合同的日计息方式累计；表中利率为合同年利率，不是90天费率。


\textbf{应付利息与本金。} 计入应付利息不等于当日把利息加入贷款本金。对本次直接贷款，一般暂缓期间未付利息不因该暂缓结束自动资本化；非补贴贷款延期结束时，尚未支付的应计利息适用计入本金的处理。本次会按服务机构结算安排支付全部新增利息，不留作后续未付余额。


\textbf{办理阶段。} 两边都在正常月供到期之前提交，批准期间一致。一般暂缓申请仍由服务机构审核；本次已确认批准条件，不把它视为无条件自动生效。


\textbf{比较范围。} 本次选择上述两种短期暂停机制，评价的是90天期间新增的借款人应付利息。原还款计划、两笔贷款的其他权利及恢复付款安排保持不变，本次不办理收入驱动还款计划、减免或再融资。

\end{document}